% Verbatim copy of /Users/alingalatan/workspaces/grothe/handoff/manuscript/legacy/ch28-coherent-wick-a.tex (file dated 2026-10-25). Only this line was added.
% Consolidated source module (formerly Chapter 28). Owner: lower-bound/09-one-driver.tex
\section{Exact Conjecture K for coherent Wick A-couplings and a dimension-free finite-coupling extension}
\label{ch:coherent-wick-a}
\chapsummary{Exact coherent one-driver A comparison, unequal-sum neighborhoods and profile transport. This proof remains an input to the complete A/B theorem in \autoref{ch:complete-coherent-driver}; its earlier coverage limitations are local to this contribution.}

\paragraph{Provenance and status (IDE note, 20 October).} Lower-bound GPT's \code{proof/chapters/ch28-coherent-wick-a.tex} from the package
\pathref{gpt-reviews/coherent-wick-a-20261020/grothendieck-coherent-wick-a-20261020.zip}, ported verbatim except for this header and the closing reading note. \emph{Our replay} \st{C}:
\code{python3 -I verify\_all.py} (Python 3.12, mpmath 1.3.0, NumPy 2.3.5, Apple clang 17) from an unrelated directory: \textbf{PASS FULL EXACT-K PACKAGE in 771\,s}
(4{,}538 numerical boxes, 11{,}189 vertices, 994 analytic leaves, 106 tail leaves, 5{,}637 splits; strict numerical clearance $1152257/2^{46}$; analytic constants, Hessian, Rice envelope,
Wick/Fourier constants, finite unequal-$A$ completion, transport and the independent tree audit all PASS); log \pathref{gpt-reviews/coherent-wick-a-20261020/replay\_ide.log}.


This chapter concerns a genuinely shared driver. It proves exact
Conjecture K for the complete one-driver coherent Wick A-family,
with no amplitude or phase restriction. An explicit analytic neighborhood
joins a certified numerical exterior. The proof also gives a finite-coupling
theorem for unequal Wick sums in arbitrary dimension, and transports the
one-driver result to the supplied profile rectangle with slack less than
$10^{-4}$. These are special-class results, not a complete proof of
Conjecture K for arbitrary cubics or a new Grothendieck
endpoint. The current repository manuscript records the separate bound
$K_G^{\C}>1.401674103\stale{22 Oct}{now $\KG>1.401695322458$}$; none of the statements below changes that endpoint.

\subsection{Definitions and the exact all-amplitude result}
Let $Z,X,G$ be independent standard proper complex Gaussians, normalized
by $\E|Z|^2=1$. Set
\[
 T=|Z|^2,\qquad H(X)=X(|X|^2-2)/\sqrt2,
 \qquad Q=cH(X)+\alpha(T-1)X+gG.
\]
Here $c,g\ge0$ and $\alpha\in\C$. Arbitrary complex coefficients of
$H$ and $G$ reduce to this form by rotations of their Gaussian drivers.
For $L_{a,b}=Z(a+bT)$ define
\[
 \Lambda_{a,b,\varepsilon}(Q)
 =\E\bigl(|L_{a,b}+Q|+|L_{a,b}-Q|\bigr)
    -\frac{\E|Q|^2}{\varepsilon},
\]
\[
 S(a,b,\varepsilon)
 =\sup_{v\ge0}\left(2\E|L_{a,b}+\sqrt vG|
                  -v/\varepsilon\right).
\]
The price is exactly
$\E|Q|^2=c^2+|\alpha|^2+g^2$; in particular the coherent sum is not
replaced by an independent sum. Write
\[
 (a_0,b_0,\varepsilon_0)
   =(699/1250,179/1250,1891/2500),\qquad
 S_0=S(a_0,b_0,\varepsilon_0),
 \quad \eta_0=\frac{703312310}{2^{46}}<10^{-5}.
\]

\begin{theorem}[Exact coherent Wick A comparison]
\label{thm:coherent-a-exact}
For every $c,g\ge0$ and $\alpha\in\C$,
\[
 \boxed{\Lambda_{a_0,b_0,\varepsilon_0}
       \bigl(cH(X)+\alpha(T-1)X+gG\bigr)\le S_0.}
\]
The supremum over this complete family equals $S_0$. Every maximizer
has $c=\alpha=0$ and $g^2$ equal to the unique maximizing Gaussian
variance.
\end{theorem}
The analytic neighborhood is proved in
Theorem~\ref{thm:coherent-sharp}; the exhaustive exterior certificate
and the proof of the exact conclusion are described below. We first
record the independently replayed numerical comparison from which the
exact cover was refined.

\begin{theorem}[A purely numerical comparison, retained as a checkpoint]
\label{thm:coherent-a-global}
For every $c,g\ge0$ and $\alpha\in\C$,
\[
 \Lambda_{a_0,b_0,\varepsilon_0}(Q)
 \le \frac{96684070713622}{2^{46}}
 \le S_0+\eta_0<S_0+10^{-5}.
\]
\end{theorem}

\paragraph{Exact radial reduction.}
Let $R=|X|^2$ and let $\theta$ be uniform on $[0,2\pi]$, independent
of $T,R$. Put
\[
 x=\sqrt T(a_0+b_0T),\quad
 m=\sqrt R\bigl(c(R-2)/\sqrt2+\alpha(T-1)\bigr),
 \quad F_w(u)=\E|\sqrt u+\sqrt wG|.
\]
The two signs in the payoff have equal integrals, and the exact norm
term is $2\E F_{g^2}(|x+me^{i\theta}|^2)$. This formula retains the
relative phase between the Wick and centered A terms. Conditionally
rotate the uniform angle by $\arg m$ (arbitrarily when $m=0$).
Since $F_w$ is
increasing and concave in $u$, its integral on an angular cell with
midpoint $\theta_j$ and half-width $d$ is at most
\[
 F_w\left(x^2+|m|^2+2x|m|\frac{\sin d}{d}\cos\theta_j\right).
\]
This is conditional Jensen applied to the exact squared displacement.
There is no pointwise optimization of the relative phase and no
ellipse-ordering assumption.

\paragraph{Radial integration and complete parameter cover.}
In the radii $r=\sqrt T,s=\sqrt R$, all affine inputs to the norm are
polynomials of degree at most three in each variable. Positive tensor
Bernstein integration bounds the norm by the weighted norm of its
controls. The complex control $cH_k+\alpha A_jX_k$ is formed before
taking its modulus. The exact norm term being bounded is convex in
$(c,\Re\alpha,\Im\alpha,g)$. On a parameter box of side lengths $h_i$,
\[
 \sup_{\text{box}}\Lambda
 \le\max_{\text{vertices}}\Lambda
          +\frac1{4\varepsilon_0}\sum_i h_i^2.
\]
The quadratic correction follows directly by interpolating the
negative sum-of-squares price. It does not require the numerical
majorant itself to be convex.

Conjugation permits $\Im\alpha\ge0$. If $V=c^2+|\alpha|^2+g^2<1$,
the parameters lie in $[0,1]\times[-1,1]\times[0,1]\times[0,1]$.
If $V\ge1$, Cauchy--Schwarz gives
\[
 \Lambda\le 2\sqrt{M+V}-V/\varepsilon_0<4/3,
 \qquad M=a_0^2+4a_0b_0+6b_0^2=1181331/1562500.
\]
The last expression decreases on $V\ge1$. Thus the unbounded parameter
tail is covered. For $V<1$, discarding either radius above six loses
at most $4e^{-18}<2^{-23}$ in the two-sign norm; every oracle call
pays this allowance.

\paragraph{Certificate and replay obligations.}
The fixed binary tree has 5,259 accepted boxes and 106 variance-tail
boxes, separated by 5,364 splits; 12,353 distinct vertices are evaluated.
All acceptance comparisons use outward integer arithmetic with scale
$2^{-46}$. The universal Rice barrier, its interpolation bounds, and
positive radial and angular data are independently checked and regenerated.
The Gaussian lower benchmark is
\[
 B=96683367401312/2^{46}\le S_0,
\]
obtained at variance $1/25$ by Jensen in visible amplitude on 4,096
radial cells. Subtracting $B$ from the achieved upper bound gives
exactly $\eta_0$. The accompanying fixed-tree checker never refines the
tree. Its runner reconstructs the integration data, checks every leaf
and both unbounded tails, and verifies $100000(703312310)<2^{46}$.
The written reduction together with these checks proves
Theorem~\ref{thm:coherent-a-global}; this is not a Lean formalization.

\subsection{Profile transport and conditional mixtures}
\begin{lemma}[Transport of an excess over a Gaussian benchmark]
\label{lem:coherent-transport}
For square-integrable visible terms $L,L_0$ on the same probability
space and any admissible $Q$,
\[
 |\Lambda_{L,\varepsilon}(Q)-\Lambda_{L_0,\varepsilon}(Q)|
 \le2\E|L-L_0|.
\]
Their Gaussian benchmarks satisfy the same bound. Consequently a
uniform comparison $\Lambda_{L_0,\varepsilon}(Q)\le S_{L_0,\varepsilon}+\eta$
transports to
\[
 \Lambda_{L,\varepsilon}(Q)\le S_{L,\varepsilon}
                         +\eta+4\E|L-L_0|.
\]
No independence between $Q$ and the visible coordinate is required.
\end{lemma}
\begin{proof}
Apply the reverse triangle inequality to each of the two norms. Their
prices cancel. Apply the same bound for each Gaussian variance and
then take the supremum to compare the benchmarks. Adding the two
inequalities proves the conclusion.
\end{proof}

\begin{corollary}[The complete supplied profile rectangle]
\label{cor:coherent-profile}
For $a\in[0.55919,0.55921]$, $b\in[0.14319,0.14321]$, and
$\varepsilon=\varepsilon_0$, all coherent Wick A-couplings satisfy
\[
 \boxed{\Lambda_{a,b,\varepsilon_0}(Q)
       < S(a,b,\varepsilon_0)+\frac{1773}{20000000}
       =S(a,b,\varepsilon_0)+0.00008865.}
\]
Relative to the fixed peak benchmark one also has
$\Lambda_{a,b,\varepsilon_0}(Q)<S_0+0.000044325$.
\end{corollary}
\begin{proof}
Gaussian radial moments give $\E|Z|=\sqrt\pi/2$ and
$\E|Z|^3=3\sqrt\pi/4$. Hence
\[
 4\E|L_{a,b}-L_{a_0,b_0}|
 \le2\sqrt\pi|a-a_0|+3\sqrt\pi|b-b_0|
 <\frac{1773}{20000000},
\]
using $\sqrt\pi<1773/1000$ and the zero-slack peak theorem.
When the reference benchmark stays $S_0$, only the first Lipschitz
bound is needed, halving the transport cost.
\end{proof}

\begin{corollary}[Simultaneous rescaling of the profile and price]
For any $\lambda>0$, put $\varepsilon=\lambda\varepsilon_0$.
For all members of the same coherent family,
\[
 \Lambda_{a,b,\varepsilon}(Q)
 \le S(a,b,\varepsilon)
     +2\sqrt\pi|a-\lambda a_0|+3\sqrt\pi|b-\lambda b_0|.
\]
\end{corollary}
\begin{proof}
The exact scaling identity is
$\Lambda_{\lambda L_0,\lambda\varepsilon_0}(Q)
=\lambda\Lambda_{L_0,\varepsilon_0}(Q/\lambda)$; the Gaussian
benchmark scales in the same way. The family is closed under division
by $\lambda$. Now use Lemma~\ref{lem:coherent-transport}.
\end{proof}

The exact all-amplitude and transported bounds also hold for
\[
 Q=C H(X)+A(T-1)X+\sqrt V G
\]
when $(C,A,V)$ is jointly random, independent of $(Z,X,G)$, $V\ge0$,
and $\E(|C|^2+|A|^2+V)<\infty$. Condition on the coefficients and
integrate the deterministic comparison. This permits dependence
among the remote coefficients, but not dependence of those coefficients
on $X$ or $Z$. It does not identify an arbitrary cubic buffer with a
Gaussian variance mixture.

\begin{corollary}[Improper Gaussian buffers and conditional covariance]
\label{cor:coherent-improper}
Let $\mathcal F$ be independent of $(Z,X)$. Let $C,A$ and a positive
semidefinite real $2$-by-$2$ matrix $\Sigma$ be $\mathcal F$-measurable.
Conditionally on $\mathcal F$, let $B$ be a centered complex Gaussian,
independent of $(Z,X)$, with real covariance $\Sigma$. If
$\E(|C|^2+|A|^2+\operatorname{tr}\Sigma)<\infty$, then
\[
 \Lambda_{a_0,b_0,\varepsilon_0}
       \bigl(CH(X)+A(T-1)X+B\bigr)\le S_0.
\]
This includes a real Gaussian buffer. Equality requires $C=A=0$ and
$\Sigma=(v_*/2)I$ almost surely, where $v_*$ is the unique Gaussian
maximizing variance.
\end{corollary}
\begin{proof}
Condition on $\mathcal F$. Simultaneously rotating $Z$ and $X$ shows
that each of $L\pm(CH(X)+A(T-1)X)$ has a circular law. Its expected
norm after adding an independent buffer therefore depends only on
the buffer's radius. An independent uniform phase may be applied to
$B$ without changing either norm or its variance.

Let $\lambda_1,\lambda_2\ge0$ be the eigenvalues of $\Sigma$ and
$\Theta$ an independent uniform angle. Rotational averaging of the
Gaussian characteristic function shows that the phase-randomized
buffer has the law $\sqrt V G$, where
\[
 V=2(\lambda_1\cos^2\Theta+\lambda_2\sin^2\Theta),
 \qquad \E[V\mid\mathcal F]=\operatorname{tr}\Sigma.
\]
Condition on $V$ and apply Theorem~\ref{thm:coherent-a-exact}. Averaging
preserves the quadratic price exactly. Equality in the averaged
comparison requires $C=A=0$ and $V=v_*$ almost surely. The latter
forces $\lambda_1=\lambda_2=v_*/2$.
\end{proof}
The covariance may depend on independent latents, but not on the
coherent driver $X$ or visible variable $Z$. Conditional nonzero
means and arbitrary non-Gaussian buffers are not covered by this
corollary. The transported profile bounds above apply to this class
as well.

\begin{proposition}[Quantitative loss of improper Gaussian noise]
For a centered complex Gaussian $B$ independent of $Z$, put
$v=\E|B|^2$, $p=|\E B^2|$, $M=\E|L|^2$, and
$f(v)=2\E|L+\sqrt vG|$. At the peak profile,
\[
 \Lambda(B)\le f(v)-v/\varepsilon_0
       -\frac{\sqrt\pi\,p^2}{16(M+v+p)^{3/2}}.
\]
This additional penalty is asserted only for the pure-buffer case.
\end{proposition}
\begin{proof}
The preceding phase randomization represents the buffer by a proper
variance mixture $V=v+p\cos(2\Theta)$, with mean $v$, variance $p^2/2$,
and range $[v-p,v+p]$. The visible mixture formula gives
$f(w)=\sqrt\pi\,\E\sqrt{R_0+w}$ with $\E R_0=M$. Jensen implies
\[
 -f''(w)\ge\frac{\sqrt\pi}{4(M+w)^{3/2}}
          \ge\frac{\sqrt\pi}{4(M+v+p)^{3/2}}
\]
throughout the mixing interval away from zero; continuity extends
the corresponding strong concavity to zero. Its Jensen deficit is
at least half this constant times $\operatorname{Var}(V)=p^2/2$.
Since $\Lambda(B)=\E f(V)-v/\varepsilon_0$, the result follows.
\end{proof}

\subsection{An explicit exact-K neighborhood}
\label{sec:coherent-local}
The small-amplitude conclusion is stronger than a slack bound.
\begin{theorem}[Exact comparison near the Gaussian buffer]
\label{thm:coherent-local}
At the exact peak profile, if
\[
 \frac7{200}\le g^2\le\frac9{200},\qquad
 0\le c\le\frac1{250},\qquad |\alpha|\le\frac1{65},
\]
then $\Lambda(Q)\le S_0-(2/5)c^4$.
On the smaller interval $0\le c\le1/425$, with the other ranges
unchanged, one has
\[
 \boxed{\Lambda(Q)\le S_0-2c^4-|\alpha|^2/8.}
\]
Both conclusions allow every complex relative phase.
\end{theorem}

\paragraph{The exact profile yields a fourth-order deficit.}
The positive Gaussian scale-mixture representation for the marginal
law of $L$ gives $L\overset d=\sqrt R G_0$, with
$\E R=M$. Use this representation only in the remote-only payoff,
never for the joint pair $(L,T)$ in an A term. With
$f(v)=2\E|L+\sqrt vG|=\sqrt\pi\E\sqrt{R+v}$,
\[
 -f''(v)>3/5,\qquad
 \frac{f'''(v)}{-f''(v)}\le\frac3{2v}
 \quad (7/200\le v\le9/200).
\]
The first inequality is Jensen using $M+v<81/100$.
For a fixed shift, $P_v(u)=2\E|L+\sqrt vG+\sqrt u|$ is a Bernstein
function. Its derivatives at zero satisfy
$P_v^{(n)}(0)=f^{(n)}(v)/n!$, so the global cubic majorant is
\[
 P_v(u)\le f(v)+u f'(v)+u^2f''(v)/4+u^3f'''(v)/36.
\]
The Wick moments are $1,44,11352$ at orders two, four, six. Comparing
the resulting expectation to $f(v+c^2)$ gives
\[
 \Lambda(cH+gG)
 \le S_0-[-f''(g^2)]\left(\frac{21}{2}
                       -\frac{473c^2}{g^2}\right)c^4
 \le S_0-\frac{63}{10}c^4+\frac{56760}{7}c^6
\]
for $c\le1/40$. The supplied mixture proof and Laplace formula justify
the derivative signs; a fixed shift is not replaced by added Gaussian
noise.

\paragraph{The coherent content starts at cubic order.}
Set $\psi_g(z)=\E|z+gG|$ and $J_c(\alpha)=\Lambda(Q)$.
Its full fourth derivative has operator norm at most
$3\sqrt\pi/(4g^3)<204$ on this buffer interval. This follows by
differentiating the absolutely convergent Fourier integral for the
smoothed norm. In $D_\alpha J_c(0)$ the terms of degrees zero, one,
and two in $c$ vanish: respectively by centering, Wick orthogonality,
and the circular phase. Thus
\[
 \|\nabla_\alpha J_c(0)\|
 \le\frac{204}{3}c^3\E|T-1|\E(|H|^3|X|)<600c^3.
\]
Here $\E|T-1|=2/e<3/4$ and
$\E(|H|^3|X|)=(-32+480e^{-2})/(2\sqrt2)<35/3$.

\paragraph{Uniform curvature and square completion.}
The independently replayed directed-interval Hessian certificate gives
$\tfrac12D_h^2J_0(0)<-7/20$ for every complex unit direction $h$ and
every buffer variance in the interval. Its checked integrand is
\[
 2\E[D_g(T)(T-1)^2]-1/\varepsilon_0,\qquad
 D_g(T)=\frac{\sqrt\pi}{4g}e^{-x}I_0(x),\quad
 x=\frac{T(a_0+b_0T)^2}{2g^2}.
\]
Taylor's formula for the Hessian, circular cancellation of its linear
term, and the fourth derivative bound give
\[
 \tfrac12D_h^2J_c(\alpha)
 \le-\frac7{20}
       +102\bigl(4c^2+4\sqrt2c\Re\alpha+18|\alpha|^2\bigr).
\]
Integrating along $t\alpha$ retains the Taylor weight, and with
$r=|\alpha|$ yields the useful master inequality
\begin{align*}
 \Lambda(Q)\le S_0-\frac{63}{10}c^4+\frac{56760}{7}c^6
 +600c^3r
 -\left[\frac7{20}-102(4c^2+2cr+3r^2)\right]r^2.
\end{align*}
On $c\le1/250$, $r\le1/65$, the bracket exceeds $1/4$.
Use $600c^3r-r^2/4\le360000c^6$ and the exact comparison
\[
 \frac{63}{10}
 -\frac{56760/7+360000}{250^2}
 =\frac{17949}{43750}>\frac25.
\]
To retain the A penalty, use
$600c^3r-r^2/8\le720000c^6$ and
\[
 \frac{63}{10}
 -\frac{56760/7+720000}{425^2}
 =\frac{1147521}{505750}>2.
\]
This proves Theorem~\ref{thm:coherent-local}. The accompanying expanded
proof gives the Gaussian Fourier identities and every Wick moment;
the rational checker verifies these constants. The Hessian replay
covers the whole variance interval and the infinite radial tail.

\subsection{Averaging the visible profile first: stronger content and curvature}
The preceding explicit neighborhood can be enlarged substantially. This
section retains the visible average before bounding derivatives, in both
the signed A profile and the A-squared profile. In the notation above,
put $r=|\alpha|$ and keep $g^2\in[7/200,9/200]$.

\begin{theorem}[A stronger quantitative coherent comparison]
\label{thm:coherent-sharp}
For every $c,r\ge0$ the following global polynomial upper bound holds:
\begin{align*}
 \Lambda(Q)\le S_0-21c^4+\frac{20812}{3}c^6
 +76c^3|\Re\alpha|
 -\left(\frac7{20}-48c^2-357cr-306r^2\right)r^2.
\end{align*}
In particular, $\Lambda(Q)\le S_0-2c^4$ throughout each of the rectangles
\[
 (c,r)\in[0,1/30]\times[0,1/100],\quad
 [0,1/32]\times[0,1/80],\quad
 [0,1/36]\times[0,1/65].
\]
One also has $\Lambda(Q)\le S_0-6c^4-r^2/50$ for
$c\le1/40$, $r\le1/65$.
\end{theorem}

\paragraph{Signed visible densities and radial convolution.}
Set $h(t)=t(a_0+b_0t)^2$ and $d(t)=h'(t)$. For $j=1,2$ define
\[
 F_{j,g}(x)=\E[(T-1)^j|L+x+gG|]=K_{j,g}(|x|^2).
\]
Before Gaussian convolution its source density on the plane is
$k_j(|x|^2)/\pi$, where
\[
 k_j(h(t))=(t-1)^j e^{-t}/d(t).
\]
Distributionally $\Delta|x|=1/|x|$ in two dimensions; there is no atom
at zero, because the flux across a circle of radius $\rho$ tends to
zero with $\rho$. Therefore
\[
 \Delta^2F_{j,0}=(\Delta(k_j(|\cdot|^2)/\pi))*|\cdot|^{-1}.
\]
For a radial nonincreasing majorant $\bar m_j(\rho^2)$ of
$4|k_j'(\rho^2)+\rho^2k_j''(\rho^2)|$, layer-cake integration of
intersections of balls gives
\[
 \|\Delta^2F_{j,0}\|_\infty
 \le2\int_0^\infty\bar m_j(\rho^2)\,d\rho.
\]
Gaussian convolution cannot enlarge this bound. The standalone
rational checker constructs 800 cells for each weight and includes
the infinite tails, obtaining respectively
\[
 \|\Delta^2F_{1,g}\|_\infty<52,\qquad
 \|\Delta^2F_{2,g}\|_\infty<64.
\]
For a radial function $F(x)=K(u)$, $u=|x|^2$,
$\Delta^2F=16(u^2K'')''$. Twice integrating from zero proves
\[
 \|K_{1,g}''\|_\infty<13/8,\qquad
 \|K_{2,g}''\|_\infty<2.
\]
This avoids taking a separate large derivative supremum for every
visible Gaussian value. The source polynomials, positive decreasing
majorants, and tail formulas are printed in the expanded proof.

\paragraph{The content bound, with its phase retained.}
With $R=|X|^2$, the exact directional derivative is
\[
 D_\zeta J_c(0)
 =2\sqrt2 c\Re\zeta\,
   \E\left[R(R-2)K_{1,g}'\bigl(c^2R(R-2)^2/2\bigr)\right].
\]
Since $\E[R(R-2)]=0$, subtract $K_{1,g}'(0)$ and use its Lipschitz
bound. The result is
\[
 |D_\zeta J_c(0)|
 \le \sqrt2\frac{13}{8}(-32+480e^{-2})c^3|\Re\zeta|
 <76c^3|\Re\zeta|.
\]
The derivative in the imaginary A direction is exactly zero. This is
a bound on the true coherent content, not a replacement by an
independent-driver calculation.

\paragraph{Inverse moments strengthen the pure residual deficit.}
For the visible mixing variable $R_0$ let $Y=1/R_0$. Its first three
moments, computed from the visible radial density at zero, are
\[
 m_1=a_0^{-2},\quad m_2=a_0^{-4}+4b_0a_0^{-5},\quad
 m_3=a_0^{-6}+12b_0a_0^{-7}+42b_0^2a_0^{-8}.
\]
H\"older applied to $Y=[Y/(1+vY)](1+vY)$ gives
\[
 \E\left(\frac{Y}{1+vY}\right)^{3/2}
 \ge\frac{m_1^{3/2}}{\sqrt{1+3vm_1+3v^2m_2+v^3m_3}}.
\]
Together with the maximum
$\sup_{y\ge0}y^{1/2}/(1+vy)^{5/2}=16/(25\sqrt{5v})$,
exact rational comparisons imply $-f''(v)>2$ and $f'''(v)<22$ on
the buffer interval. Substitution in the Bernstein majorant proves,
for every $c\ge0$,
\[
 \Lambda(cH+\sqrt vG)\le S_0-21c^4+(20812/3)c^6.
\]

\paragraph{The A-squared average strengthens the curvature.}
The radial Hessian identity gives
\[
 \|D^2F_{2,g}(x)-D^2F_{2,g}(0)\|
 \le6\|K_{2,g}''\|_\infty |x|^2<12|x|^2.
\]
Using $\E(|X|^2|H|^2)=4$, the A curvature at $\alpha=0$ is at most
$-7/20+48c^2$, replacing the old cost $408c^2$.
For every complex unit direction $\zeta$, expansion in $\alpha$ then gives
\[
 \tfrac12D_\zeta^2J_c(\alpha)
 \le-7/20+48c^2+1071cr+1836r^2.
\]
Here the linear term at $c=0$ vanishes by the X phase; the fourth
derivative bound $204$ and the moments
$\E|T-1|^3<5/2$, $\E(|H||X|^3)<21/10$ give $1071$.
The quadratic term uses $\E(T-1)^4=9$ and $\E|X|^4=2$.
Integrating along the A segment divides the last two coefficients
by three and six respectively. This proves the polynomial bound in
Theorem~\ref{thm:coherent-sharp}.

For the three rectangles, replace the bracket by its value at their
upper-right corner and complete the square in $r$. The exact
remaining quartic coefficients are respectively
\[
 \frac{1773616}{744525}>2,\qquad
 \frac{2339375}{1138944}>2,\qquad
 \frac{58173389}{19509012}>2.
\]
Retaining $r^2/50$ before completing the square proves the last
assertion. For a purely imaginary A coefficient, the linear term
vanishes, giving the further region
\[
 c\le1/20,\quad |\alpha|\le1/100,\quad\Re\alpha=0
 \quad\Longrightarrow\quad
 \Lambda(Q)\le S_0-3c^4-|\alpha|^2/50.
\]
All rectangle comparisons are replayed with exact fractions.

\begin{corollary}[Dimension-free first variation for unequal Wick sums]
\label{cor:coherent-diffuse-content}
Let $X_1,\ldots,X_n$ be independent proper complex Gaussians, and let
$B$ be an independent circular square-integrable random variable.
For $q=\sum_j c_jH(X_j)+B$ and
$Q_{\boldsymbol\alpha}=q+(T-1)\sum_j\alpha_jX_j$,
\[
 \left\|\nabla_{\boldsymbol\alpha}
          \Lambda(Q_{\boldsymbol\alpha})\big|_{\boldsymbol\alpha=0}
 \right\|_2
 \le76\left(\sum_j|c_j|^6\right)^{1/2}.
\]
In particular, coefficients $c_j=c/\sqrt n$ give $76|c|^3/n$.
\end{corollary}
\begin{proof}
Convolution by any circular probability law preserves radiality and
the bound $\|\Delta^2F_1\|_\infty<52$. Thus the signed-visible content
proof works with any independent circular square-integrable buffer,
not just a Gaussian buffer. For coordinate $j$, regard the remaining
summands and $B$ as that buffer; the corresponding gradient block
has Euclidean norm at most $76|c_j|^3$. Sum the squares of these
blocks. The price has zero first derivative because of Wick
orthogonality and centering. The smooth visible density justifies
the derivatives also for a nonsmooth buffer law, by convolution and
dominated differentiation.
\end{proof}

\subsection{Closing the all-amplitude coherent family}
The exact cover uses the same verified norm oracle, parameter root,
radial tails, and variance tail as the purely numerical checkpoint.
Its leaf types have different mathematical conclusions and must not
be conflated:
\begin{itemize}
\item A numerical leaf bounds its exact payoff by $B\le S_0$.
\item An analytic leaf is contained in one of the three rectangles of
Theorem~\ref{thm:coherent-sharp}, including the whole buffer-variance
interval condition. It gives $\Lambda\le S_0-2c^4$ directly.
\item A variance-tail leaf gives $\Lambda<4/3<B\le S_0$.
\end{itemize}
The exact binary cover contains 4,538 numerical leaves, 994 analytic
leaves, and 106 variance-tail leaves, with 5,637 splits. The checker
verifies the analytic membership conditions using the maximum squared
complex A modulus on each box, not the modulus of its midpoint. It
also checks the minimum and maximum squared buffer amplitudes. All
tree descendants, the absence of trailing records, and all tail
conditions are checked from the original root box.

The fixed-tree replay regenerates the integration data, reruns the
analytic scalar and Hessian certificates, and recompiles its C++
checker with assertions enabled. Its receipt separates the maximum
\emph{numerical exterior} upper bound from the analytic bound $S_0$.
That numerical maximum is not an upper bound on the full family,
because the analytic leaves contain the Gaussian maximizer.

Every parameter belongs to one of these leaves or to the unbounded
variance tail. Therefore $\Lambda\le S_0$ everywhere, proving the
inequality in Theorem~\ref{thm:coherent-a-exact}. Conversely the family
contains every $gG$, so its supremum is exactly $S_0$.

The numerical leaves have strictly positive clearance below $B$;
the replay reports their exact rational maximum. On an analytic leaf,
the sharper local polynomial bound is strictly below $S_0$ if $c>0$.
If $c=0$ and $\alpha\ne0$, it gives
\[
 \Lambda(Q)\le S_0-\bigl(7/20-306|\alpha|^2\bigr)|\alpha|^2<S_0.
\]
Thus every maximizer is Gaussian. Finally
$f(v)=\sqrt\pi\E\sqrt{R_0+v}$ is strictly concave, and
$f(v)-v/\varepsilon_0\to-\infty$ as $v\to\infty$; its maximizing
variance is unique. The cover places it in $[7/200,9/200]$.

This argument closes the exact one-driver A conjecture by combining
analytic estimates with a complete exterior computation. It does not
deduce a global zero-slack comparison from the earlier small numerical
slack alone.

\begin{corollary}[Uniform stability within the coherent family]
For every $c,g\ge0$ and $\alpha\in\C$ at the peak profile,
\[
 S_0-\Lambda(Q)\ge
 \min\left\{10^{-10},\,2c^4+\frac{|\alpha|^2}{10000}\right\}.
\]
\end{corollary}
\begin{proof}
On each analytic rectangle reserve $|\alpha|^2/10000$ in the sharp
polynomial inequality before completing the square. The three retained
quartic coefficients are
\[
 \frac{7067548}{2976075},\qquad
 \frac{58134775}{28449024},\qquad
 \frac{11565279133}{3897367164},
\]
all greater than $2$. Thus the second term in the minimum is a lower
bound on every analytic leaf. The complete fixed-tree replay gives
maximum numerical exterior upper numerator $96683366249055$, so its
clearance below $B$ is
\[
 \frac{1152257}{2^{46}}>10^{-10}.
\]
The variance tail has $B-4/3>10^{-10}$. These comparisons, including
the retained quartic coefficients, are checked with exact fractions
by the independent partition audit.
\end{proof}

\subsection{Finite shared A-couplings in arbitrary dimension}
The preceding first-variation bound can also pay for finite couplings,
with no bound on the number of unequal Wick blocks.
\begin{theorem}[Unequal Wick sums with simultaneous coherent A]
\label{thm:coherent-unequal}
Let $X_1,\ldots,X_n,Z,G$ be independent standard proper complex
Gaussians, let $c_j,\alpha_j\in\C$, and set
\[
 q=\sum_jc_jH(X_j),\quad
 Q=q+\sqrt vG+(T-1)\sum_j\alpha_jX_j,\quad
 B_4=\sum_j|c_j|^4.
\]
At the exact peak profile, assume
\[
 \frac7{200}\le v\le\frac9{200},\qquad
 u:=\sum_j|c_j|^2\le\frac1{100},\qquad
 c_*:=\max_j|c_j|\le\frac1{50},\qquad
 r:=\left(\sum_j|\alpha_j|^2\right)^{1/2}\le\frac1{1000}.
\]
Then
\[
 \boxed{\Lambda(Q)\le S_0-2B_4-r^2/50.}
\]
\end{theorem}

\paragraph{A sharper buffered Fourier deficit.}
The inherited characteristic-function certificate proves
$\phi(s)=\E e^{is\Re H(X)}\ge e^{-s^2/4}$ for $0\le s\le4$.
The included replay checks this continuous interval with outward
interval arithmetic, verified quadrature nodes, and analytic tails.
An elementary stronger estimate, using the sixth-order cosine lower
bound and the quadratic exponential upper bound, gives
\[
 \phi(s)\ge\exp(-s^2/4+(3/5)s^4),\qquad 0\le s\le1/10.
\]
Indeed the difference of those polynomial bounds divided by $s^4$,
with $t=s^2$, is
$9/160-(473/96-3/20)t-9t^2/50>0$ on $[0,1/100]$.

The visible characteristic function $\psi(k)$ is positive by its
Gaussian-mixture representation and satisfies
$\psi(k)\ge e^{-Mk^2/4}$. The exact norm Fourier identity, together
with independence, gives a low-frequency deficit at least
\[
 \frac65 B_4 I,\qquad
 I=\int_0^5 k^2e^{-(M+11/200)k^2/4}\,dk>4.7686.
\]
The lower bound is a rational integral of the degree-25 lower Taylor
polynomial for the exponential. Frequencies between $5$ and
$K=4/c_*$ give nonnegative deficit. Beyond $K$, retaining the
Gaussian buffer bounds the possibly negative contribution by
\[
 \frac4K e^{-vK^2/4}
 =c_*e^{-4v/c_*^2}<\frac{B_4}{1000}.
\]
For the last inequality use $B_4\ge c_*^4$, $v\ge7/200$, and
monotonicity of $x^{-3}e^{-7/(50x^2)}$ on $x\le1/50$; its maximum
is $125000e^{-350}<1/1000$. The case $B_4=0$ is Gaussian and is
handled directly. The resulting pure-residual estimate is
\[
 \Lambda(q+\sqrt vG)\le S_0-\frac{17}{3}B_4.
\]

\paragraph{Paying for all shared directions at once.}
For a unit vector $h$, write $P_h=\sum_jh_jX_j$. Wick orthogonality
gives the exact mixed moment
\[
 \E(|q|^2|P_h|^2)=u+3\sum_j|c_j|^2|h_j|^2
 \le D:=u+3c_*^2.
\]
The averaged A-squared estimate therefore gives A curvature at zero
at most $-7/20+12D$. In its variation along the coupling segment,
weighted Cauchy--Schwarz gives
$\E(|q||P_h|^3)\le\sqrt{2D}$. The same fourth-derivative argument
as before proves
\[
 \Lambda(Q)\le\Lambda(q+\sqrt vG)
 +76\sqrt{B_6}\,r
 -\left(\frac7{20}-12D-170\sqrt{2D}\,r-306r^2\right)r^2,
 \quad B_6=\sum_j|c_j|^6.
\]
The bracket exceeds $9/50$ on the displayed parameter set. Reserve
$r^2/50$ and complete the square with the remaining $4r^2/25$:
\[
 76\sqrt{B_6}\,r-\frac4{25}r^2\le9025B_6
 \le\frac{361}{100}B_4.
\]
Since $17/3-361/100=617/300>2$, the theorem follows.

For a concrete application-range example, take 25 coefficients
$c_j=1/50$ and Gaussian variance $v=7/200$. The remote residual has
normalized fourth moment
\[
 2+\frac{42B_4}{(u+v)^2}=2+\frac{56}{675}
 =2.08296296\ldots.
\]
At $v=1/25$ the value is $2.0672$. Thus this finite-coupling theorem
enters the lower end of the application's kurtosis window. Its
restriction $\|\boldsymbol\alpha\|_2\le10^{-3}$ remains essential to
the present proof; it is not an all-amplitude or all-cubic result.

\subsection{What this resolves and what still blocks the endpoint}
The coherent family violates the conditional phase symmetries used in
the preceding shared-sector comparisons. The exact theorem therefore
adds a different case; it does not follow by declaring the Wick residual
conditionally Gaussian. A weak coherent-A comparison was already
available with larger slack. It is now superseded by an exact comparison
at the peak, with a complete all-amplitude cover.

The arbitrary-dimensional theorem also covers finite coherent A
couplings on unequal Wick sums, including examples in the requested
kurtosis window, but with explicit coefficient and coupling restrictions.
All-amplitude A/B couplings for the coherent one-driver family are
subsequently closed in \autoref{ch:complete-coherent-driver}; the applicable
even-sector extensions are in \autoref{ch:all-sector-audit-extension}.
Arbitrary cubic self-interaction and unrestricted simultaneous
visible-sector interactions remain outside those theorems.
The profile-rectangle error below $10^{-4}$ applies only to the
one-driver A family. A complete admissible row envelope and a comparison
for all remaining residuals are still required before changing
$K_G^{\C}$. In particular, the numerical B phase relaxation must not be
reused as if it were the honest coherent B payoff: its artificial
quadratic gain obstructs an exact comparison before any numerical
resolution issue arises.

\subsection{IDE reading note}\label{sec:coherent-reading}
\paragraph{Independent check \st{F}.} Our own code (\pathref{forest-2026-10-20/coherent\_a\_own\_check.py}, not GPT's) evaluates $\Lambda$ by quadrature after the joint phase rotation
($T$ Gauss--Laguerre $120$, $X$ polar $120\times96$, the buffer integrated in closed form by the Rice function), optimising $g$ for each of $190$ points
$c\in\{0,\dots,0.5\}$, $|\alpha|\in\{0,\dots,0.3\}$ and four phases. Result: $\max(\Lambda-S)=-2\cdot10^{-16}$ (no violation), and $S-\Lambda\ge2c^4$ at every tested point of the analytic
neighbourhood ($c\le1/30$, $|\alpha|\le1/100$, $g^2\in[.035,.045]$). Example: $S-\Lambda=2.33\cdot10^{-7}$ at $(c,\alpha)=(0.01,0)$ against $2c^4=2\cdot10^{-8}$. The phase of $\alpha$
matters little ($\le15\%$ of $S-\Lambda$), and $\alpha=0.005$ already costs $1.16\cdot10^{-5}$, the profile loss of \autoref{ch:forest}.
\paragraph{What it settles.} \autoref{ch:kjoint}'s Wick-junta $A$-only bound $S+7.3\cdot10^{-5}$ is now exact ($\eta=0$), so this family uses none of the budget $\eta_{AB}$. Still open:
$B$ with a shared driver; interactions with $Y,W,\nu$; several shared drivers at finite amplitude (the dimension-free theorem of this chapter needs $\sum c_j^2\le0.01$ and $\|\alpha\|\le10^{-3}$,
i.e.\ $\sigma_q\le0.1$, far from $\sigma^*\approx0.2$). The method, averaging the visible coordinate before bounding derivatives, is the principle of \autoref{ch:forest}~(a): it is the
reason this proof closed where the pointwise constants ($600c^3$, $408c^2$) could not.
